\chapter*{MỞ ĐẦU}
\addcontentsline{toc}{chapter}{MỞ ĐẦU}
\markboth{MỞ ĐẦU}{}

\section*{1. Tính cấp thiết của đề tài}
\addcontentsline{toc}{section}{1. Tính cấp thiết của đề tài}

Các hệ thống thông tin vô tuyến thế hệ mới phải đồng thời phục vụ số lượng thiết bị ngày
càng lớn, duy trì tốc độ truyền dữ liệu cao, bảo đảm chất lượng dịch vụ và rút ngắn thời
gian ra quyết định. Trong bối cảnh phổ tần và công suất phát luôn hữu hạn, việc chỉ tăng
công suất không còn là giải pháp đủ hiệu quả. Hiệu năng hệ thống phụ thuộc ngày càng
nhiều vào khả năng quản lý nhiễu giữa các người dùng và khả năng điều khiển môi trường
truyền sóng.

Khái niệm đa truy cập (Multiple Access, MA) mô tả cách nhiều người dùng cùng khai thác
tài nguyên vô tuyến. Các kỹ thuật truyền thống như đa truy cập trực giao (Orthogonal
Multiple Access, OMA), đa truy cập phân chia theo không gian (Space Division Multiple
Access, SDMA) và đa truy cập phi trực giao (Non-Orthogonal Multiple Access, NOMA) giải
quyết bài toán chia sẻ tài nguyên theo các nguyên lý khác nhau. Tuy nhiên, khi kênh
truyền, mức nhiễu và yêu cầu của người dùng thay đổi, không có một cơ chế duy nhất luôn
tối ưu trong mọi tình huống. Đa truy cập phân chia tỷ lệ (Rate-Splitting Multiple Access,
RSMA) được đề xuất như một cách quản lý nhiễu linh hoạt hơn bằng việc tách thông điệp
thành phần chung và phần riêng, qua đó cho phép người nhận vừa giải mã một phần nhiễu vừa
xem phần còn lại như tạp âm \cite{Mao2018RSMABridging, Mao2022RSMAFundamentals}.

Song song với việc thay đổi cách tổ chức tín hiệu ở phía phát, bề mặt phản xạ thông minh
có thể cấu hình lại (Reconfigurable Intelligent Surface, RIS) cung cấp khả năng điều khiển
môi trường truyền sóng bằng nhiều phần tử thụ động có thể thay đổi pha. RIS truyền thống
chủ yếu phản xạ tín hiệu về một nửa không gian. Bề mặt đồng thời truyền và phản xạ
(Simultaneously Transmitting and Reflecting RIS, STAR--RIS) mở rộng khả năng đó bằng việc
đồng thời tạo thành phần truyền qua và phản xạ, nhờ đó có thể phục vụ người dùng ở cả hai
phía của bề mặt \cite{Xu2021STARRIS, Mu2021STARAided}.

Khi STAR--RIS kết hợp với RSMA, số biến điều khiển tăng nhanh. Bộ điều khiển không chỉ
phải chọn công suất cho luồng chung và các luồng riêng mà còn phải quyết định cách chia
tốc độ luồng chung cho từng người dùng, cách chia năng lượng truyền và phản xạ trên từng
phần tử STAR--RIS cùng pha tương ứng của hai phía. Các biến này tác động lẫn nhau thông
qua kênh hiệu dụng và các biểu thức tỷ số tín hiệu trên nhiễu cộng tạp âm
(Signal-to-Interference-plus-Noise Ratio, SINR). Vì vậy, bài toán tối ưu có tính phi lồi,
số chiều cao và phải được giải lại khi thông tin trạng thái kênh (Channel State
Information, CSI) thay đổi.

Các phương pháp tối ưu lặp như tối ưu luân phiên (Alternating Optimization, AO) có thể
tìm được nghiệm cục bộ có chất lượng cao, nhưng chi phí tính toán tăng mạnh khi số phần tử
STAR--RIS tăng. Trong môi trường kênh biến đổi nhanh, việc chạy lại hàng trăm hoặc hàng
nghìn phép đánh giá hàm mục tiêu cho mỗi CSI có thể trở thành nút thắt về độ trễ.

Một hướng tiếp cận khác là sử dụng học tăng cường sâu (Deep Reinforcement Learning, DRL)
để học trước ánh xạ từ CSI sang quyết định tài nguyên. Sau giai đoạn huấn luyện, mạng
nơ-ron chỉ cần một lượt truyền xuôi (forward pass) để tạo hành động, nhờ đó giảm đáng kể
thời gian tính toán trực tuyến. Tuy nhiên, việc dùng DRL không tự động bảo đảm hiệu quả.
Nếu không gian hành động được thiết kế quá tổng quát, mạng phải tự học từ dữ liệu cả cấu
trúc vật lý của bài toán lẫn giá trị quyết định tối ưu. Đối với hệ thống một đầu vào, một
đầu ra (Single-Input Single-Output, SISO), vấn đề này đặc biệt rõ vì không có tiền mã hóa
theo không gian để tách các luồng riêng. Nhiều luồng riêng cùng hoạt động có thể gây nhiễu
lẫn nhau mạnh, nhất là trong miền tỷ số tín hiệu trên tạp âm cao. Phân tích này gợi ý rằng
việc đưa kiến thức vật lý vào cách biểu diễn hành động có thể quan trọng không kém việc
lựa chọn thuật toán DRL.

Từ những lý do trên, luận văn nghiên cứu bài toán phân bổ tài nguyên trong hệ thống SISO
STAR--RIS hỗ trợ RSMA theo hướng kết hợp hai lớp: lớp mô hình vật lý xác định cấu trúc
hợp lý của hành động và lớp DRL học phần điều chỉnh còn lại theo CSI. Ba thuật toán DRL
được khảo sát gồm Deep Deterministic Policy Gradient (DDPG), Twin Delayed Deep
Deterministic Policy Gradient (TD3) và Proximal Policy Optimization (PPO). Các thuật toán
này được đối chiếu với các phương pháp tối ưu lặp và giải tích để đánh giá đồng thời tổng
tốc độ, khả năng đáp ứng chất lượng dịch vụ (Quality of Service, QoS), độ ổn định và độ
trễ ra quyết định.

\section*{2. Vấn đề nghiên cứu}
\addcontentsline{toc}{section}{2. Vấn đề nghiên cứu}

Trong hệ thống được xét, một trạm gốc (Base Station, BS) SISO truyền dữ liệu đồng thời cho
bốn thiết bị người dùng (User Equipment, UE) thông qua một STAR--RIS thụ động. RSMA tạo
một luồng chung được tất cả UE giải mã và bốn luồng riêng tương ứng với bốn UE. Sau khi
giải mã luồng chung, UE thực hiện khử nhiễu liên tiếp (Successive Interference
Cancellation, SIC) để loại luồng chung rồi mới giải mã luồng riêng của mình.

Một CSI cụ thể dẫn đến một bài toán quyết định gồm năm nhóm biến vật lý: công suất luồng
chung và luồng riêng; tỷ lệ phân bổ tốc độ chung; hệ số chia năng lượng truyền và phản xạ;
pha truyền; và pha phản xạ. Mục tiêu là tăng tổng tốc độ của hệ thống trong khi duy trì
tốc độ tối thiểu của từng UE. Hai yêu cầu này có thể mâu thuẫn: dồn tài nguyên cho UE có
kênh tốt làm tổng tốc độ tăng nhanh nhưng có thể khiến UE yếu vi phạm QoS; chia đều tài
nguyên giúp công bằng hơn nhưng có thể làm tăng nhiễu giữa các luồng riêng và giảm hiệu
quả sử dụng công suất.

Một câu hỏi quan trọng của luận văn là liệu mạng DRL có nên tự học toàn bộ miền hành động
vật lý hay nên nhận trước một số cấu trúc có căn cứ từ mô hình. Mã nguồn thực nghiệm cho
thấy cách biểu diễn hành động có cấu trúc gồm ba ý chính: tỷ lệ công suất luồng chung được
giới hạn trong một vùng ưu tiên; phần công suất luồng riêng được chia bằng một hàm softmax
có độ sắc (sharpness) cao để có khả năng tạo phân bổ gần một luồng chiếm ưu thế; và pha STAR--RIS được
học như phần hiệu chỉnh quanh một pha tham chiếu giải tích phụ thuộc kênh. Các lựa chọn
này không phải định luật vật lý bắt buộc, mà là các giả thiết thiết kế nhằm thu hẹp miền
tìm kiếm về vùng có ý nghĩa đối với bài toán đang xét.

Vì vậy, vấn đề nghiên cứu không được đặt theo hướng chứng minh một thuật toán DRL cụ thể
luôn vượt trội. Thay vào đó, luận văn trả lời ba câu hỏi chính. Thứ nhất, cấu trúc vật lý
của SISO--RSMA gợi ý điều gì về cách phân bổ công suất và pha? Thứ hai, khi dùng cùng một
bộ giải mã hành động vật lý, nhóm thuật toán học ngoài chính sách (off-policy) gồm DDPG và
TD3 khác nhóm học theo chính sách hiện hành (on-policy) là PPO như thế nào? Thứ ba, nhóm
DRL tạo ra đánh đổi giữa chất lượng nghiệm và độ trễ ra sao so với nhóm tối ưu lặp AO?

\section*{3. Mục tiêu nghiên cứu}
\addcontentsline{toc}{section}{3. Mục tiêu nghiên cứu}

\subsection*{3.1. Mục tiêu tổng quát}

Mục tiêu tổng quát của luận văn là xây dựng và đánh giá một khung DRL có biểu diễn hành
động dựa trên cấu trúc vật lý cho bài toán phân bổ tài nguyên trong hệ thống SISO
STAR--RIS hỗ trợ RSMA, qua đó làm rõ đánh đổi giữa chất lượng nghiệm, khả năng đáp ứng QoS
và độ trễ ra quyết định.

\subsection*{3.2. Mục tiêu cụ thể}

Các mục tiêu cụ thể gồm:

\begin{enumerate}[leftmargin=*,itemsep=2pt]
\item Xây dựng mô hình kênh, tín hiệu RSMA, hệ số STAR--RIS và bài toán tối ưu nhất quán
      với phần mềm mô phỏng.
\item Phân tích đặc điểm nhiễu của các luồng riêng trong hệ SISO để giải thích vì sao phân
      bổ công suất tập trung và tỷ lệ công suất luồng chung cao có thể tạo vùng hành động
      thuận lợi trong cấu hình nghiên cứu.
\item Thiết kế bộ giải mã hành động có cấu trúc để biến đầu ra chuẩn hóa của mạng nơ-ron
      thành các biến vật lý khả thi gồm công suất, phân bổ tốc độ chung, hệ số chia năng
      lượng và pha STAR--RIS.
\item Triển khai và đánh giá ba thuật toán DRL gồm DDPG, TD3 và PPO trên cùng mô hình vật
      lý, cùng không gian hành động và cùng giao thức dữ liệu.
\item Xây dựng các phương pháp đối chiếu gồm AnalyticalRIS, AO--SCA và AO--Grid; so sánh
      nhóm DRL với nhóm tối ưu lặp về tổng tốc độ, QoS và độ trễ tính toán.
\item Thiết kế giao thức thực nghiệm có tập huấn luyện, xác thực và kiểm thử tách biệt, sử
      dụng cùng các kịch bản kênh khi so sánh phương pháp và lưu thông tin phục vụ tái lập
      kết quả.
\end{enumerate}

\section*{4. Đối tượng và phạm vi nghiên cứu}
\addcontentsline{toc}{section}{4. Đối tượng và phạm vi nghiên cứu}

Đối tượng nghiên cứu thứ nhất là hệ thống truyền xuống SISO STAR--RIS hỗ trợ RSMA với một
BS, một STAR--RIS thụ động và bốn UE. Đối tượng thứ hai là bài toán phân bổ công suất và
tốc độ luồng chung trong RSMA. Đối tượng thứ ba là bài toán điều khiển hệ số chia năng
lượng và pha của STAR--RIS. Đối tượng thứ tư là các thuật toán DRL liên tục và các bộ giải
tối ưu lặp được dùng để tạo quyết định tài nguyên theo CSI.

Phạm vi nghiên cứu được giới hạn như sau. Kênh được sinh tổng hợp theo phân bố Gaussian
phức với các hệ số tỷ lệ cố định trong mã nguồn; CSI được giả định biết đầy đủ tại thời
điểm ra quyết định; SIC được giả định lý tưởng; STAR--RIS hoạt động theo cơ chế chia năng
lượng (Energy Splitting, ES); số UE cố định bằng bốn; số phần tử STAR--RIS thay đổi trong
tập $\{16, 32, 64, 96, 128\}$.

Luận văn không xét sai số ước lượng CSI, méo phần cứng, lượng tử hóa pha, nhiều anten tại
BS, mô hình pha truyền và phản xạ ghép phần cứng hoặc triển khai trên thiết bị vô tuyến
thực. Các kết luận thực nghiệm vì vậy chỉ có ý nghĩa trong phạm vi mô hình và giao thức đã
nêu. Đặc biệt, các giá trị như khoảng công suất luồng chung, độ sắc của softmax hoặc biên
hiệu chỉnh pha là lựa chọn thiết kế cho cấu hình đang nghiên cứu, không được diễn giải như
hằng số tối ưu chung cho mọi hệ thống STAR--RIS--RSMA.

\section*{5. Phương pháp nghiên cứu}
\addcontentsline{toc}{section}{5. Phương pháp nghiên cứu}

Luận văn kết hợp nhiều phương pháp. Phương pháp mô hình hóa toán học được dùng để xây dựng
kênh hiệu dụng, biểu thức SINR, tốc độ và ràng buộc. Phương pháp phân tích xấp xỉ được
dùng để giải thích hành vi của nhiễu luồng riêng trong miền tỷ số tín hiệu trên tạp âm cao
và từ đó hình thành trực giác cho biểu diễn hành động. Phương pháp học tăng cường sâu được
dùng để huấn luyện DDPG, TD3 và PPO. Phương pháp tối ưu số được dùng để xây dựng AO--SCA,
AO--Grid và AnalyticalRIS làm đối chiếu. Cuối cùng, phương pháp mô phỏng Monte Carlo và
phân tích thống kê được dùng để đánh giá trên các kịch bản kênh độc lập.

Một nguyên tắc xuyên suốt là mọi công thức, biến, cấu hình mạng và chỉ số đánh giá trong
luận văn phải truy ngược được về mã nguồn hoặc tài liệu tham khảo. Khi một nhận định chỉ
là trực giác, xấp xỉ hoặc đề xuất mở rộng, nội dung đó được ghi rõ thay vì trình bày như
kết quả đã được chứng minh.

\section*{6. Ý nghĩa khoa học và thực tiễn}
\addcontentsline{toc}{section}{6. Ý nghĩa khoa học và thực tiễn}

Về mặt khoa học, luận văn nhấn mạnh vai trò của biểu diễn hành động trong DRL cho bài toán
tối ưu vô tuyến. Thay vì xem mạng nơ-ron như một hộp đen phải tự khám phá toàn bộ hình học
của miền khả thi, nghiên cứu đưa một số cấu trúc từ mô hình SISO--RSMA và căn chỉnh pha
STAR--RIS vào bộ giải mã. Cách tiếp cận này tạo cầu nối giữa phân tích vật lý và học máy:
mô hình vật lý xác định vùng tìm kiếm hợp lý, còn mạng nơ-ron học sự thích nghi theo CSI
trong vùng đó.

Về mặt thực tiễn, kết quả cung cấp một đánh giá rõ về đánh đổi giữa hai nhóm phương pháp.
Các bộ giải AO có thể đạt tổng tốc độ cao hơn vì thực hiện tối ưu lặp trực tiếp trên mỗi
CSI và tìm kiếm trên miền vật lý rộng. Ngược lại, DRL sau huấn luyện chỉ cần một lượt
truyền xuôi qua mạng nơ-ron nên có độ trễ ra quyết định thấp hơn nhiều. Kết quả này giúp
định vị DRL như một lựa chọn phù hợp cho các tình huống kênh thay đổi nhanh và yêu cầu
quyết định trực tuyến, thay vì xem DRL là sự thay thế tuyệt đối cho tối ưu số.

\section*{7. Đóng góp của luận văn}
\addcontentsline{toc}{section}{7. Đóng góp của luận văn}

Các đóng góp chính của luận văn gồm bốn nhóm. Thứ nhất, xây dựng một mô hình tính toán
thống nhất cho STAR--RIS, RSMA và QoS để các thuật toán học và tối ưu sử dụng cùng một hàm
đánh giá vật lý. Thứ hai, phân tích cấu trúc nhiễu trong hệ SISO để giải thích xu hướng
phân bổ luồng riêng tập trung và vai trò của luồng chung trong việc duy trì QoS. Thứ ba,
xây dựng bộ giải mã hành động có cấu trúc gồm giới hạn tỷ lệ công suất luồng chung, softmax
có độ sắc cao cho công suất luồng riêng, phép chiếu lên đơn hình (simplex) cho phân bổ tốc độ chung,
ánh xạ năng lượng khả thi và pha tham chiếu kèm phần hiệu chỉnh học được. Thứ tư, xây dựng
giao thức thực nghiệm so sánh ba thuật toán DRL với các phương pháp tối ưu lặp trên các
tập kênh tách biệt, đồng thời báo cáo cả chất lượng nghiệm và độ trễ.

\section*{8. Kết cấu luận văn}
\addcontentsline{toc}{section}{8. Kết cấu luận văn}

Ngoài phần mở đầu, kết luận, tài liệu tham khảo và phụ lục, luận văn gồm bốn chương.
Chương~\ref{chap:co-so-ly-thuyet} trình bày nền tảng về đa truy cập, RSMA, RIS và
STAR--RIS, DRL cùng các phương pháp tối ưu truyền thống; đồng thời tổng hợp nghiên cứu
liên quan và xác định khoảng trống. Chương~\ref{chap:mo-hinh-he-thong} xây dựng mô hình hệ
thống, mô hình tín hiệu, bài toán tối ưu và phân tích cấu trúc SISO--RSMA.
Chương~\ref{chap:phuong-phap} trình bày biểu diễn trạng thái, bộ giải mã hành động có cấu
trúc, kiến trúc mạng nơ-ron, ba thuật toán DRL, các phương pháp đối chiếu và giao thức
thực nghiệm. Chương~\ref{chap:ket-qua} trình bày kết quả mô phỏng, so sánh nhóm DRL với
nhóm tối ưu lặp, đánh giá QoS, độ trễ, khả năng mở rộng và các giới hạn của nghiên cứu.
